\section{NN-S012: three different supported-zero circuit restrictions}
\label{sec:root-circuit-models}

\paragraph{The models and exact statement.}
Let $R$ be a nonzero commutative unital ring. Circuits have fan-in two
primitive gates $\oplus,\odot$ in $G(R)$, constants in $G(R)$, and
Boolean inputs encoded by $0\mapsto\tau$, $1\mapsto1_R$.
Write $e=0_R$ for the supported zero, $\pi:G(R)\to R$ for amplitude,
and $\sigma:G(R)\to\mathbb B$ for support. A circuit computes a
Boolean function by amplitude if its final amplitude is always the
element $0_R$ or $1_R$ representing that function. Distinguish the
following three restrictions:
\begin{enumerate}
\item Reading only final support $\sigma$ imposes no restriction on
      intermediate amplitudes.
\item Requiring the final value to be exactly $\tau$ or $1_R$ forbids
      a supported zero at the output, but permits it at internal gates.
\item Requiring that no gate, including a constant gate, takes the
      value $e$ on any Boolean input forbids supported zeros at every
      primitive gate.
\end{enumerate}
In the third model the final amplitude is again required to be Boolean.
Circuit sizes count primitive gates, with a fixed constant cost for each
ring constant.

For every $R$, models 1 and 2 compute precisely the monotone Boolean
functions, and their minimum circuit sizes are within constant factors
of monotone Boolean circuit size. The same holds for model 3 whenever
$R\not\cong\mathbb F_2$. Over $\mathbb F_2$, however, model 3 computes
only constants and conjunctions of subsets of the input variables.
In particular, two-variable OR has no such circuit over $G(\mathbb F_2)$,
although it has a constant-size such circuit over $G(\mathbb F_4)$.

If $R$ is fixed and finite, unrestricted amplitude computation over
$G(R)$ is equivalent, up to constant factors in circuit size, to ordinary
Boolean computation. The finite-ring hypothesis is needed for this
automatic simulation: no claim about unit-cost arbitrary integer or
complex constants is implicit.

\paragraph{Proof of the support and output statements.}
At every gate the pair $(\pi,\sigma)$ evolves by ring addition and
multiplication in the first coordinate and OR and AND in the second.
The support of each constant is fixed. Induction on the circuit therefore
gives a monotone circuit of the same size for every gate's support.
This proves the lower bounds and monotonicity in model 1. In model 2,
the support of the output equals its Boolean amplitude, proving the same
lower bound and monotonicity. In model 3 this equality also holds at the
Boolean output.

For the upper bound in model 1, replace Boolean OR by $\oplus$ and AND
by $\odot$; only the support is read. For model 2, use multiplication
for AND and the following constant-size circuit for OR on inputs
$a,b\in\{\tau,1_R\}$:
\[
 (a\oplus b)\oplus((-1_R)\odot(a\odot b)).             \tag{S012.1}
\]
Its amplitude is $a+b-ab$, with $\tau$ sent to zero. Its support is
the Boolean OR of the two inputs. Checking their four possible values
shows that its final value is $\tau$ for $(\tau,\tau)$ and $1_R$
otherwise. Thus each Boolean gate is replaced by a fixed number of
primitive gates and its output has the required encoding. Internal
supported zeros have not been excluded in this construction.

\paragraph{Proof for every-gate absence.}
First let $R=\mathbb F_2$ and suppose a circuit satisfies model 3.
Every gate takes values in $\{\tau,1\}$. If both inputs of an addition
gate have nonempty support as functions of the input bits, their
monotone support functions both take value one at the all-ones input.
The two input values there are both $1$, so the addition gate takes
the forbidden value $e=1\oplus1$. Hence at least one input of each
addition gate is identically $\tau$. Delete that gate and retain its
other input. This can also be done effectively by evaluating support
at the all-ones input: a monotone support is identically zero exactly
when it vanishes there. What remains has only multiplication gates,
the constants $\tau,1$, and input variables. It computes a constant
or a conjunction. Conversely all those functions have such circuits.

If $2\ne0$ in $R$, circuit (S012.1) already satisfies model 3.
The first sum has nonzero amplitude $1$ or $2$ whenever supported;
the product and its negative have amplitude $1$ and $-1$ whenever
supported, and the output has just been checked. All constants used
are nonzero. This proves the monotone upper bound in that case.

Suppose instead that $2=0$ and $R\not\cong\mathbb F_2$.
Choose any $\lambda\in R\setminus\{0,1\}$; no unit hypothesis is
needed. The following circuit computes OR without a supported zero
at any primitive gate:
\[
 s=a\oplus(\lambda\odot b),\qquad
 h=(1+\lambda)\oplus a,\qquad
 t=b\odot h,\qquad s\oplus t.                         \tag{S012.2}
\]
The nonzero supported values of $s$ belong to
$\{1,\lambda,1+\lambda\}$. The gate $h$ is always supported and has
value $1+\lambda$ or $\lambda$, both nonzero. The gate $t$ is absent
when $b=\tau$, and otherwise has one of these same two nonzero values.
The final value is $\tau,1,1,1$ on the ordered inputs
$(\tau,\tau),(1,\tau),(\tau,1),(1,1)$, respectively. All constants
and gates therefore satisfy the restriction. AND is multiplication
on $\{\tau,1\}$, so substituting these circuits at monotone gates
proves the upper bound for every remaining nonzero ring. The lower
bound was the support shadow. This completes the classification.

\paragraph{Proof of the finite-ring amplitude simulation.}
Ordinary AND and NOT are represented in amplitude by $ab$ and $1-a$.
Lifting them to $G(R)$ preserves these amplitude identities by the
homomorphism $\pi$, even when a false intermediate value is $e$ rather
than $\tau$. Thus every Boolean circuit has a constant-factor
arithmetic lift. Conversely, when $R$ is fixed and finite, every
element of $G(R)$ has a fixed-length binary encoding. Each addition
or multiplication is a fixed finite lookup table and has a Boolean
circuit of constant size. Encode all gates and extract the final
amplitude bit. This gives the reverse constant-factor simulation.
The elementary proof uses no E, A, M, or unproved complexity separation.

\paragraph{What is corrected in the supplied circuit dictionary.}
The source is \emph{OURS\_TRANSLATION\_SPLIT\_ZERO\_20261121.md},
Section 9, items 1--6. Its support-shadow assertion is correct.
The phrase ``no supported zeros'' must specify model 2 or model 3.
In characteristic two the displayed OR formula
$x\oplus y\oplus xy$ has an internal supported zero at $x=y=1$.
Thus its proposed monotone equivalence is correct for output-only
absence, but false for absence at every primitive gate over
$\mathbb F_2$. Passing to $\mathbb F_4$ repairs that precise gate
restriction with (S012.2). It does not prove a Boolean circuit lower
bound: unrestricted computation over either fixed finite field still
has the same Boolean circuit complexity.

For a monotone family $f_n$, let $M_n$ be its monotone circuit size
and $C_n$ its general Boolean circuit size. The exact corrected
comparison is between $M_n$ and $C_n$, or the equivalent split models
specified above. In particular, $\mathrm{CLIQUE}\in\mathrm P/\mathrm{poly}$
is equivalent to a polynomial upper bound in the unrestricted amplitude
model. A statement merely that the ratio $M_n/C_n$ is exponential
does not provide that upper bound: as numerical growth conditions,
$M_n=2^{n^a}$ and $C_n=2^{n^a/2}$ have an exponential ratio and neither
is polynomial. This example is an explanation of the missing logical
premise, not an assertion of these actual CLIQUE complexities.
Any claimed equivalence using ``exponential savings'' must include
the polynomial endpoint, or supply an additional theorem proving it.
The same caution applies to E: its proved subexponential running times
give subexponential circuits for the encoded computations, not a
polynomial endpoint. NN-C001 and NN-C019 give the exact supported-zero
removal and instance-dependent isolation consequences available from E.

\paragraph{Geometric formulation and proof dictionary.}
Over $\mathbb F_2$, put
\[
 X=\operatorname{Spec}
 \mathbb F_2[x_1,\ldots,x_n]/(x_1^2-x_1,\ldots,x_n^2-x_n).
\]
This is the finite disjoint union of the $2^n$ rational points indexed
by Boolean assignments. The amplitude at a gate is a section of its
finite free function algebra; its support is an additional upward-closed
subset of that indexing Boolean poset. In model 3 the support is exactly
the nonvanishing locus of the section. Two nonempty upward-closed
supports contain the all-ones point, and their sections both equal one
there. Their sum vanishes there in characteristic two while the support
tag remains present. This is precisely the addition obstruction proved
above. It is a statement about the actual presentation by functions and
support tags, not about a virtual Grothendieck class or a Frobenius weight.

The finite-free algebra and its base-change trace are the setting of
EGA IV$_4$, 18.2.1, especially equation (18.2.1.2), and 18.2.2,
original printed page 111. The original page image was read in this
run. Base change to $\mathbb F_4$ keeps the same split finite scheme
and supplies the additional constant $\lambda$ used in (S012.2).
Thus every-gate zero avoidance can change under a finite scalar
extension even though the Boolean function and its support poset do
not change. The support tags are extra data beyond this algebra;
NN-S003 explains why additive ringization forgets them. For a general
ring, absence of the zero element at an $R$-valued input is not
silently identified with invertibility or with a base-change-stable
open condition. The general-ring theorem was proved directly in
$G(R)$, and the displayed scheme dictionary is specifically the
finite-field case. Neither purity nor A or M imposes the missing
polynomial upper or lower circuit bound.
\hfill$\square$
